% Author biography for the ASEJ submission portal (Attach Files -> Author biography).
% Corresponding author only, by the authors' decision. If the portal or the editorial
% office asks for one per author, the remaining five skeletons are in
% author_biographies.md and the questionnaire is author_bio_form.md.
%
% This file identifies the author by design, like the title page. It must never be
% merged into manuscript.tex, which is anonymized for double-blind review.
\documentclass[12pt,a4paper]{article}
\usepackage[T1]{fontenc}
\usepackage[utf8]{inputenc}
\usepackage{mathptmx}
\usepackage[margin=2.5cm]{geometry}
\usepackage{microtype}
\pagestyle{empty}
\setlength{\parindent}{0pt}
\setlength{\parskip}{8pt}

\begin{document}

{\large\bfseries Author biography\par}

\textbf{Manuscript title:} Predictive fidelity versus training utility in surrogate-based
reinforcement learning for HVAC control: a role-separated hybrid remedy

\vspace{10pt}
\textbf{Almaz Sapargali} (corresponding author)

Almaz Sapargali received the B.Sc.\ degree in cyberphysical systems from the International
Information Technology University (IITU), Almaty, Kazakhstan, in 2023, and is currently pursuing
the M.Sc.\ degree in computer science and technology at Al-Farabi Kazakh National University,
Almaty, Kazakhstan, under a double-degree programme with Northwestern Polytechnical University
(NPU), China. Research interests include deep reinforcement learning and scientific machine
learning, with a focus on surrogate-based control of building energy systems. Earlier results from
the present study were presented at the WCCM--ECCOMAS Congress 2026 (Munich, Germany, 19--24 July
2026) under the title ``Deep reinforcement learning for energy-efficient control in smart
buildings''. In the present work the contributions covered the methodology, the software, the
experiments and the original draft.

\end{document}
